% ---------------------------------------------------------------------------------
% 3. RESULTS AND ERROR ANALYSIS
% Guide: macro-F1 + per-class P/R/F1 (train / val / public apart); experiment table;
% curves + confusion matrices; visual error analysis; public scores with submission IDs.
% ---------------------------------------------------------------------------------
\section{Results and error analysis}\label{sec:results}

\textbf{Metrics.} The main metric is macro-F1 (mean of the seven per-class F1), so rare classes count as much as frequent ones. Validation F1 drove every decision; training F1 is only a diagnostic; the public CodaBench score was never used for selection.

\textbf{What helped and what did not} (\cref{tab:experiments}):
\begin{itemize}
  \item \textbf{CUSTOM:} the reference reached $\sim$0.36 and our first GAP version 0.440, both underfitting. The search plateaued at 0.45--0.465. More capacity (48 channels, 2 convs/block) and 60 epochs gave the selected model: \textbf{0.502}.
  \item \textbf{No gain:} inverse-frequency weights (worse), 64 channels, 5 blocks, larger crops, CutMix, 80 epochs, ensembles, per-class threshold tuning and retraining on train+val. Seeds alone move results by $\pm0.008$.
  \item \textbf{FT:} every modern backbone beat AlexNet (0.452); Swin-T was best with \textbf{0.546}.
\end{itemize}

\textbf{Curves and confusion matrices.} For CUSTOM (\cref{fig:curves}) train and validation curves stay together and validation F1 saturates at 0.49--0.50 after epoch 30: no overfitting, the model is limited by the data. Swin-T fits the training set better (train F1 0.77 vs 0.546 on validation). \Cref{fig:cm} shows the validation confusion matrices (row-normalised \%) \todo{describe once generated}.

\textbf{Error analysis} (\cref{tab:perclass}, \cref{fig:errors}).
\begin{itemize}
  \item \textbf{Easy:} Water (0.85) and Non-burnable (0.80).
  \item \textbf{Hard:} the intermediate levels Low (0.22), Moderate (0.32) and High (0.33) for CUSTOM, and Low, Moderate and Very High for FT, which points to confusions between neighbouring risk levels \todo{confirm with \cref{fig:cm,fig:errors}}.
  \item \textbf{Limited signal:} even on its own training images (no augmentation) CUSTOM only reaches 0.53 (Low 0.30), so the problem is the information in the image, not model capacity.
  \item \textbf{CUSTOM vs FT:} CUSTOM is small (2.7\,M parameters), fast and stable; FT gains $+0.04$ on validation at a higher compute cost.
\end{itemize}

\textbf{Competition} (\cref{tab:submissions}). Best public scores: \textbf{0.49} CUSTOM (p4\_09\_B, submission 1036062) and \textbf{0.48} FT (swin\_p2\_1). Every model loses about 0.06 from validation to public, because the same validation images were reused for hundreds of trials and the public subset is small. For example, the 121-trial model had the best validation score (0.524) but only 0.47 on public.
